hola

hem de fer:

Objetivo:
● Entender cómo encontrar estructura y patrones en datos sin etiquetar mediante
algoritmos de clustering y detección de anomalías.
Guión:
● Algoritmos de clustering: K-Means (sklearn.cluster.KMeans), DBSCAN
(sklearn.cluster.DBSCAN) y clustering jerárquico.
● Modelos de mezcla gaussiana: sklearn.mixture.GaussianMixture.
● Detección de anomalías: Isolation Forest (sklearn.ensemble.IsolationForest).
● Cómo evaluar un clustering sin etiquetas: silhouette score y método del codo.
● Implementar cuatro problemas con aplicación física. Algunas sugerencias:
o Fácil: agrupar estrellas por temperatura y luminosidad con K-Means y
comparar los grupos con el diagrama de Hertzsprung-Russell.
o Media: identificar cúmulos estelares en datos reales del satélite Gaia mediante
DBSCAN (librería astroquery).
o Difícil: detectar tránsitos de exoplanetas como anomalías en curvas de luz del
telescopio Kepler (librería lightkurve).
